% profiles/new-york-university.tex — New York University
% ============================================================================
% source: https://gsas.nyu.edu/content/dam/nyu-as/gsas/documents/dissertationsubmissionrelated/NYU%20GSAS%20Dissertation%20Formatting%20Guidelines%202025.pdf
%         (Graduate School of Arts and Science, "Dissertation Formatting
%         Guidelines", 2025)
% accessed: 2026-09-02
% checked:  top-20 research-institution pass
%
% SCOPE: these are the GSAS guidelines. NYU's other schools (Steinhardt, Tandon,
% Wagner, Silver) publish their own, and a candidate outside GSAS should not
% assume this profile applies.
%
% VERIFIED and set here: geometry, spacing.
% LEFT TO THE NEUTRAL FALLBACK: title page and approval page -- see the note at
% the bottom, which matters here more than at most schools in this set.

\thesisinstitution{New York University}

% "Margins: 1" on all sides." The guidelines repeat it against the specimen
% ("Note that the table does not extend into the margins").
\thesissetgeometry{left=1in, right=1in, top=1in, bottom=1in}

% "Text: at least 10pt legible font, double spaced." Footnotes are single-spaced
% but must still be at least 10pt. Lists, "including the Table of Contents, must
% be double-spaced", with subsections within lists "at least 1.5 spaced".
\thesissetspacing{double}

% --- Accessibility -----------------------------------------------------------
% NYU's font note is advice about on-screen reading rather than a conformance
% requirement, and is recorded as such.
\thesisaccessibilityrequirement{On-screen readability}
	{The guidelines flag certain faces as web fonts "designed for easy screen
	 readability. Since many readers are likely to view and/or use your
	 dissertation" onscreen, using one may improve readability. Advice, not a
	 requirement.}
\thesisaccessibilityrequirement{Typeface}
	{"Typeface/fonts must be at least 10 pt. or higher for the entire
	 dissertation including footnotes."}
\thesisaccessibilityrequirement{Not published}
	{The 2025 GSAS guidelines state no assistive-technology requirement: no
	 tagging, alt-text or WCAG obligation is given.}

% --- Title page ---------------------------------------------------------------
% NOT overridden, and worth flagging: NYU is the one school in this set that
% puts a SIGNATURE LINE on the title page itself -- "Provide a signature line for
% your dissertation advisor above their typed" name, with stacked signature
% lines where there is more than one. The class's neutral title page has no
% signature line, so a GSAS candidate must supply the page rather than rely on
% the fallback. The approval page is therefore also left alone rather than
% suppressed.
